\answer{ex:17:1}{$0.2$~Wに対して$81$~\textmu W。生データなら$50$~pJ/bit以下が必要。}
\answer{ex:17:2}{$c=0.1$では$5.6$、$8.7$、$9.2$~bit/s、極限$9.3$~bit/s。$c=0$、$N=1000$では$65$~bit/s。}
\answer{ex:17:3}{$B=\log_2N+P\log_2P+(1-P)\log_2\frac{1-P}{N-1}=4.87$、$4.37$、$3.29$~bit/文字、すなわち$7.3$、$6.6$、$4.9$~bit/s。$10$~bit/sには毎秒$2.3$文字が必要。}
\answer{ex:17:4}{成分あたりの誤差分散は$2b/(Nm^2T)+\sigma_\xi^2$。両者の寄与が等しいのは$N\approx90$。極限は窓あたり成分あたり$\tfrac12\log_2(1+0.25/0.04)\approx1.4$~bit。}
\answer{ex:17:5}{組は$\eta_bd^2+\eta_db^2<2$、つまりおよそ$\eta<1$で収束する。$0.8$では収束、$1.1$では周期$2$の持続振動、$1.5$では非周期的、$2$では発散。学習者の速度は分けるべきである。}
\answer{ex:17:6}{閾値は$\approx4.4\sigma$。$102$~kbit/s、$0.10$~mWで、生データの$205$~mWの$2000$分の1。飽和する（$3$スパイクを超える）のは$5.7\cdot10^{-5}$のサンプルだけ。}
\answer{ex:17:7}{自由課題。$\text{SNR}\approx0.9$で約$46$~bit/s、独立ノイズで約$330$~bit/s。信号に似たノイズが原因なら、取り出せる情報はチャネル数を増やすと飽和する。符号の無知が原因なら、より良いデコーダ（たとえばスパイクのタイミングを使うもの）で増える。}
